\documentclass[12pt, a4paper]{scrartcl}
\newcommand{\CommonPath}{../../Common}
\newcommand{\coursename}{Intro to Deep Learning}
\usepackage{\CommonPath/notespreamble}

% Title
\title{Exercise Session 1\\
\large Prerequisites Self-Assessment}
\author{Magda Gregorová\\
{\small \href{mailto:magda.gregorova@thws.de}{magda.gregorova@thws.de}}}
\date{{\small March 2026 \\[0.3em]
Licensed according to
\href{https://creativecommons.org/licenses/by-sa/4.0}{CC-BY-SA 4.0}}}

\begin{document}

\maketitle

\begin{center}
\fbox{\parbox{0.85\textwidth}{
\textbf{Note:} This is not a test and will not be graded. Work through each block independently, then we will discuss solutions together. Annotate your sheet as you go — mark what felt easy, shaky, or unfamiliar. This is your personal study map for the first weeks of the course.
}}
\end{center}

\bigskip

\tableofcontents
\newpage

%----------------------------------------------------------------------
\section{Mathematical Notation}
%----------------------------------------------------------------------

\begin{enumerate}

    \item Compute $\displaystyle\sum_{i=1}^{4} x_i$ and $\displaystyle\prod_{i=1}^{4} x_i$
    for $\mathbf{x} = (2, 5, 1, 3)^\top$.

\item Translate each expression into plain English:

\begin{itemize}
    \item $w \in \mathbb{R}$
    \item $\mathbf{x} \in \mathbb{R}^{784}$
    \item $f: \mathbb{R}^n \to \mathbb{R}$
    \item $\forall i \in \{1, \ldots, n\}$
    \item $A \in \mathbb{R}^{m \times n}$
\end{itemize}

\item Read the following definition carefully:

\begin{quote}
Let $x_i^{(k)}$ denote the $i$-th feature of the $k$-th data sample, where $i
\in \{1, \ldots, d\}$ and $k \in \{1, \ldots, n\}$. The full dataset is
collected into a matrix $X \in \mathbb{R}^{n \times d}$, where $X_{ki} =
x_i^{(k)}$.
\end{quote}

\begin{enumerate}

    \item What does $x_3^{(1)}$ refer to? Where does it appear in $X$?

    \item What is the shape of $X$ if you have 200 samples and 8 features?
    Which dimension is which?

    \item A classmate writes $x_k^{(i)}$ when they mean ``the $i$-th feature of
    the $k$-th sample.'' What is wrong, and why does it matter?

    \item Write an expression for the mean of feature $i$ across all samples
    using $\sum$ notation. How many such means are there in total?

\end{enumerate}


    \item If $f(x) = x^2$ and $g(x) = x + 1$, compute $(f \circ g)(3)$ and
    $(g \circ f)(3)$.

    \item Let $h: \mathbb{R}^n \to \mathbb{R}$ and $g: \mathbb{R} \to
    \mathbb{R}$. What are the input and output spaces of $g \circ h$? What
    about $h \circ g$?

    \item A neural network layer is written as $f(\mathbf{x}) = \sigma(W\mathbf{x}
    + \mathbf{b})$ where $W \in \mathbb{R}^{m \times n}$, $\mathbf{b} \in
    \mathbb{R}^m$, $\mathbf{x} \in \mathbb{R}^n$, and $\sigma: \mathbb{R} \to
    \mathbb{R}$ is applied elementwise. What are the input and output
    dimensionalities of $f$? For a second layer $g(\mathbf{y}) =
    \sigma(V\mathbf{y} + \mathbf{c})$ to be composable as $g \circ f$, what
    constraints follow on $V$ and $\mathbf{c}$?

    \item Consider a three-layer network $NN = f_3 \circ f_2 \circ f_1$ with
    $f_1: \mathbb{R}^{784} \to \mathbb{R}^{128}$, $f_2: \mathbb{R}^{128} \to
    \mathbb{R}^{64}$, $f_3: \mathbb{R}^{64} \to \mathbb{R}^{10}$. What are
    the input and output spaces of $NN$?

\end{enumerate}

%----------------------------------------------------------------------
\section{Linear Algebra}
%----------------------------------------------------------------------

\begin{enumerate}

    \item Write the following expression in summation notation: ``the sum of all
    pairwise products $x_i y_i$ for $i = 1, \ldots, n$.'' Does this expression
    look familiar? Is there any relation to an $\ell_2$ norm, and if so what?

    \item A simple loss function used in regression is mean squared error
    $\mathrm{MSE} = \|f(\mathbf{x}) - \mathbf{y}\|_2^2$. Let $f(\mathbf{x}) =
    2\mathbf{x}$, $\mathbf{x} = (1, 2, 3)^\top$, and $\mathbf{y} = (1.5, 4.5,
    7)^\top$. Compute the MSE. What does MSE with value zero mean?

\item Let $A = \begin{pmatrix} 1 & 2 & 0 \\ 3 & 1 & 4 \end{pmatrix}$ and
$\mathbf{x} = \begin{pmatrix} 1 \\ -1 \\ 2 \end{pmatrix}$.

\begin{enumerate}

    \item Compute $A\mathbf{x}$. What is the shape of the result?

    \item Compute $AA^\top$ and $A^\top A$.

\end{enumerate}

\item Reason about dimensions:

\begin{enumerate}

    \item Let $W_1 \in \mathbb{R}^{4 \times 3}$, $W_2 \in \mathbb{R}^{2 \times
    4}$, $\mathbf{x} \in \mathbb{R}^3$. What is the shape of $W_2 W_1
    \mathbf{x}$? What about $W_1^\top W_2^\top$?

    \item Let $X \in \mathbb{R}^{n \times d}$, $W \in \mathbb{R}^{d \times m}$,
    $\mathbf{b} \in \mathbb{R}^m$. What is the shape of $XW + \mathbf{b}$?
    What needs to hold for the addition to be valid?

    \item Let $A \in \mathbb{R}^{m \times n}$, $B \in \mathbb{R}^{n \times n}$,
    $C \in \mathbb{R}^{n \times p}$, $\mathbf{v} \in \mathbb{R}^p$. What is
    the shape of $A(B + B^\top)C\mathbf{v}$? At each step, write the
    intermediate shape.

\end{enumerate}

\end{enumerate}


%----------------------------------------------------------------------
\section{Calculus \& Gradients}
%----------------------------------------------------------------------

Compute the derivative with respect to $x$:

\begin{enumerate}

    \item $f(x) = x^3 - 2x^2 + 5$

    \item $f(x) = \sigma(x) = \dfrac{1}{1 + e^{-x}}$. Show that $\sigma'(x)
    = \sigma(x)(1 - \sigma(x))$.

    \item $f(x) = \max(0, x)$ - what is the derivative?

\item Consider the following function:
\[
L(w, b) = \frac{1}{n}\sum_{i=1}^{n}\left(\sigma(wx_i + b) - y_i\right)^2,
\qquad \text{where } \sigma(z) = \frac{1}{1+e^{-z}}
\]

\begin{enumerate}

    \item Identify the composite structure. Define intermediate variables $z_i$,
    $a_i$, and $r_i$ such that $L$ can be written as a composition of simple
    functions.

    \item Using your decomposition, compute $\dfrac{\partial L}{\partial w}$
    step by step using the chain rule. Write out each intermediate partial
    derivative explicitly. You may use the result from 3.1 that $\sigma'(z) =
    \sigma(z)(1-\sigma(z))$.

    \item For $n=2$, $\mathbf{x} = (1, 2)^\top$, $\mathbf{y} = (0, 1)^\top$,
    $w = 2$, $b = 0$: evaluate $\dfrac{\partial L}{\partial w}$ numerically.
    In which direction would you adjust $w$ to decrease $L$?

\end{enumerate}

\item Consider a single linear layer with weight vector $\mathbf{w} \in \mathbb{R}^d$
and bias $b \in \mathbb{R}$, with MSE loss:
\[
L(\mathbf{w}, b) = \frac{1}{n}\sum_{i=1}^{n}
\left(\mathbf{w}^\top \mathbf{x}^{(i)} + b - y_i\right)^2
\]

Compute
    $\nabla_\mathbf{w} L$ and $\dfrac{\partial L}{\partial b}$. What is the
    shape of $\nabla_\mathbf{w} L$?

    \item Now extend to a full linear layer with weight matrix $W \in
    \mathbb{R}^{m \times d}$ and bias vector $\mathbf{b} \in \mathbb{R}^m$,
    with loss:
    \[
    L(W, \mathbf{b}) = \frac{1}{n}\sum_{i=1}^{n}
    \left\|\mathbf{y}^{(i)} - \left(W\mathbf{x}^{(i)} + \mathbf{b}\right)\right\|_2^2
    \]
    Compute $\nabla_W L$ and $\nabla_\mathbf{b} L$. What
    are their shapes? How do they relate to the shapes of $W$ and $\mathbf{b}$?

\end{enumerate}


%----------------------------------------------------------------------
\section{Python \& PyTorch Concepts}
%----------------------------------------------------------------------

\begin{enumerate}
\item What does the following return, and why?

\begin{verbatim}
x = [1, 2, 3]
y = x
y.append(4)
print(x)
\end{verbatim}

What would you change to avoid this behaviour?

\item Without running the code, predict the output shape of each line:

\begin{verbatim}
import torch
A = torch.ones(3, 4)
B = torch.ones(4, 2)
C = A @ B          # shape?
D = A.T            # shape?
E = A.sum(dim=0)   # shape?
\end{verbatim}

\end{enumerate}

%----------------------------------------------------------------------
\section{Coding \& Installation Check}
%----------------------------------------------------------------------

\begin{enumerate}

\item Using PyTorch tensors, implement the MSE loss $\mathrm{MSE} =
\frac{1}{n}\|f(\mathbf{x}) - \mathbf{y}\|_2^2$ where $f(\mathbf{x}) = 3\mathbf{x}
- 1$ for:
\[
\mathbf{x} = (1, 2, 3, 4, 5, 6, 7, 8)^\top, \qquad
\mathbf{y} = (2.1, 5.2, 7.8, 11.1, 14.9, 17.2, 20.8, 23.1)^\top
\]

\begin{enumerate}
    \item Create $\mathbf{x}$ and $\mathbf{y}$ as PyTorch tensors.
    \item Compute $\hat{\mathbf{y}} = f(\mathbf{x})$ using tensor operations.
    \item Compute the MSE manually using tensor operations - do not use any
    PyTorch loss functions.
\end{enumerate}

\item Using your symbolic result from Section~3.3, implement $\nabla_W L$ and
$\nabla_\mathbf{b} L$ manually for the following setup, with $W \in
\mathbb{R}^{2 \times 3}$, $\mathbf{b} \in \mathbb{R}^2$, $n = 4$ samples:

\[
X = \begin{pmatrix} 1 & 2 & 1 \\ 0 & 1 & 3 \\ 2 & 0 & 1 \\ 1 & 1 & 2
\end{pmatrix}, \quad
Y = \begin{pmatrix} 1 & 0 \\ 0 & 1 \\ 1 & 1 \\ 0 & 0 \end{pmatrix}, \quad
W = \begin{pmatrix} 1 & 0 & 1 \\ 0 & 1 & 0 \end{pmatrix}, \quad
\mathbf{b} = \begin{pmatrix} 0 \\ 0 \end{pmatrix}
\]

\begin{enumerate}
    \item Create all tensors in PyTorch.
    \item Compute $\nabla_W L$ and $\nabla_\mathbf{b} L$ using your symbolic
    result — no autograd.
    \item Print the shapes of $\nabla_W L$ and $\nabla_\mathbf{b} L$ and verify
    they match the shapes of $W$ and $\mathbf{b}$.
\end{enumerate}

\end{enumerate}

\section*{Acknowledgements}
This exercise sheet was prepared with the assistance of Claude, a large
language model developed by Anthropic. All content has been reviewed, edited,
and approved by the author.
\addcontentsline{toc}{section}{Acknowledgements}

\end{document}